% ----------------------------------------------------------
\begin{frame}[fragile]{ORDER BY}
  \begin{block}{Sorting Query Results}
\begin{lstlisting}[style=sqlstyle]
-- Ascending (default)
SELECT * FROM product ORDER BY price;
SELECT * FROM product ORDER BY price ASC;

-- Descending
SELECT * FROM product ORDER BY price DESC;

-- Multiple columns: primary then secondary sort
SELECT * FROM product ORDER BY brand_id, price DESC;
\end{lstlisting}
  \end{block}

  \begin{itemize}
    \item Without \texttt{ORDER BY}, the row order is \textbf{undefined} --
          the database is free to return rows in any order, and this can
          vary between executions.
    \item Multiple columns: rows are first sorted by the first column; ties
          are broken by the second column, and so on.  Each column can
          independently be \texttt{ASC} or \texttt{DESC}.
    \item \texttt{ORDER BY} is evaluated \textbf{after} \texttt{WHERE} and
          \texttt{GROUP BY} in the execution order.
    \item \texttt{NULL} values sort to the \emph{end} in MySQL by default
          (\texttt{ASC}); other DBMSs (e.g.\ PostgreSQL) sort \texttt{NULL}
          first by default.
  \end{itemize}
\end{frame}

% ----------------------------------------------------------
\begin{frame}[fragile]{GROUP BY and HAVING}
  \begin{block}{Grouping and Filtering Groups}
\begin{lstlisting}[style=sqlstyle]
-- Average price per brand
SELECT brand_id, AVG(price) AS avg_price
FROM product
GROUP BY brand_id;

-- Only brands with more than 10 products
SELECT brand_id, AVG(price), COUNT(*) AS product_count
FROM product
GROUP BY brand_id
HAVING COUNT(*) > 10;

-- Combining WHERE and HAVING
SELECT brand_id, AVG(price)
FROM product
WHERE price > 0
GROUP BY brand_id
HAVING AVG(price) > 5;
\end{lstlisting}
  \end{block}

  \begin{itemize}
    \item \texttt{GROUP BY} \textbf{collapses} all rows with the same group
          value into a single output row.
    \item Every column in \texttt{SELECT} must either appear in
          \texttt{GROUP BY} or be wrapped in an aggregate function.
    \item \texttt{WHERE} filters rows \textbf{before} grouping;
          \texttt{HAVING} filters groups \textbf{after} grouping.
  \end{itemize}
\end{frame}

% ----------------------------------------------------------
\begin{frame}[fragile]{CASE WHEN \textbf{(!)}}
  \begin{block}{Conditional Expressions in SELECT}
\begin{lstlisting}[style=sqlstyle]
SELECT name, price,
  CASE
    WHEN price > 10 THEN 'Expensive'
    WHEN price >= 5 THEN 'Medium'
    ELSE 'Cheap'
  END AS price_category
FROM product;
\end{lstlisting}
  \end{block}

  \begin{itemize}
    \item \texttt{CASE WHEN} acts like an \textbf{if-else expression}
          inside a query -- it returns different values depending on which
          condition is met.
    \item Conditions are evaluated \textbf{in order}; the first matching
          \texttt{WHEN} branch is returned and the rest are skipped.
    \item If no condition matches and there is no \texttt{ELSE},
          \texttt{NULL} is returned.
    \item \texttt{CASE WHEN} can appear in \texttt{SELECT},
          \texttt{ORDER BY}, and even inside aggregate functions
          (e.g.\ \texttt{SUM(CASE WHEN \ldots THEN 1 ELSE 0 END)}).
  \end{itemize}

  \begin{examples}{Use Case: Conditional Aggregation}
    \texttt{SUM(CASE WHEN country = 'DE' THEN 1 ELSE 0 END)} counts only
    German brands without a subquery -- a powerful pattern for pivot-style
    reports.
  \end{examples}
\end{frame}
